% TOP_PROBLEM: {"id": 2308007, "title": "Research Problems in Function Theory — Problem 8.7", "category_id": 8, "category": {"id": 8, "name": "analysis", "display_name": "Analysis", "description": "Limits, continuity, calculus, and function theory.", "slug": "analysis", "order_index": 8, "created_at": "2026-07-31T15:26:25.671Z"}, "status": "open"}
% FIRST_POSTED: null
% ENTRY_KIND: statement_only
% Imported from: batch07.tex; UnsolvedMath ID 2308007
\documentclass[11pt]{book}
\usepackage{fontspec}
\setmainfont{Latin Modern Roman}
\setsansfont{Latin Modern Sans}
\usepackage{amsmath,amssymb,mathtools}
\usepackage{unicode-math}
\setmathfont{Latin Modern Math}
\usepackage{xurl}
\usepackage[hidelinks]{hyperref}
\newcommand{\sref}[2]{\hyperlink{source:#1:#2}{[S#2]}}
\newcommand{\cataloguescope}[1]{\paragraph{Mathematical question.}#1}
\begin{document}
% UnsolvedMath ID 2308007; source catalogue code AMR-022-8007
\setcounter{section}{2308006}
\section{Research Problems in Function Theory — Problem 8.7}\label{2308007}
{\small\sffamily UnsolvedMath ID 2308007 · Analysis}
\subsection{Definitions and mathematical statement}

\paragraph{Shared notation.}$\mathbb N=\{1,2,\ldots\}$, and $\mathbb Z,\mathbb Q,\mathbb R,\mathbb C$ denote the integers, rationals, reals and complex numbers. Any other conventions are specified locally.


Let $\mathbb D=\{z\in\mathbb C:|z|<1\}$ and $\mathbb T=\partial\mathbb D$. For $p>0$, the Bergman space is $A^p(\mathbb D)=\{f\text{ analytic}:\|f\|_{A^p}=(\int_{\mathbb D}|f|^p\,dA)^{1/p}<\infty\}$, where $dA=dx\,dy$ is ordinary area measure. Let $(z_n)_{n\ge1}\subset\mathbb D$ and set $k_n(z)=(1-z_nz)^{-2}$. Suppose $M=\overline{\operatorname{span}\{k_n:n\ge1\}}^{A^2}$ is a proper subspace of $A^2$. Let $f_j$ be finite complex linear combinations of the $k_n$, with $\|f_j-f\|_{A^2}\to0$ for some $f\in A^2$. Each finite combination is a rational function; removable poles are canceled when evaluating it outside the disk. For the analogous simple-pole construction the kernels are $h_n(z)=(1-z_nz)^{-1}$, with the same finite-linear-combination and proper-closed-span conditions in a specified normed space of analytic functions. The Hardy space mentioned by the source is $H^2(\mathbb D)=\{g\text{ analytic}:\sup_{0<r<1}(2\pi)^{-1}\int_0^{2\pi}|g(re^{it})|^2\,dt<\infty\}$.
\paragraph{Statement recorded in the supplied source.}
If there exists an open disk $\Delta\subset\{|z|>1\}$ on which these rational functions are defined and $\sup_{z\in\Delta}|f_j(z)|\to0$, must $f\equiv0$? \sref{2308007}{2}
Hayman notes that similar questions can be formulated with the simple-pole kernels $h_n$ and with analytic function spaces other than $A^2$, and reports that the analogous $H^2$ result is known. \sref{2308007}{2}
\subsection{Short English statement}
Does vanishing of finite kernel combinations on an exterior disk force their Bergman-norm limit to vanish when the kernels do not span the whole space? Hayman also notes analogous questions for other kernels and spaces.
\subsection{Sources}
\begin{description}
\item[\hypertarget{source:2308007:1}{S1}] ULAM.AI and the UnsolvedMath Contributors, \emph{UnsolvedMath: A Curated Collection of Open Mathematics Problems}, record 2308007 (AMR-022-8007). CC BY 4.0. \url{https://huggingface.co/datasets/ulamai/UnsolvedMath}.
\item[\hypertarget{source:2308007:2}{S2}] Walter K. Hayman and Edward F. Lingham, \emph{Research Problems in Function Theory} (2018), arXiv:1809.07200. Problem 8.7, complete original question, all following updates, and the relevant chapter definitions. \url{https://arxiv.org/abs/1809.07200}.
\end{description}
\label{end:2308007}
\end{document}
